\documentclass[11pt]{article}

\usepackage[a4paper,margin=30mm]{geometry}
\usepackage{amsmath,amssymb,amsthm,mathtools}

\newcommand{\E}{\mathbb{E}}
\newcommand{\Pp}{\mathbb{P}}
\newcommand{\R}{\mathbb{R}}
\newcommand{\eqdef}{\mathrel{\vcentcolon=}}
\newcommand{\arginf}{\mathop{\mathrm{arg\,inf}}}

% Local notation from IMF.tex.
\newcommand{\rhohu}{\rho}
\newcommand{\rhohum}{\phi}
\newcommand{\Kern}{K}
\newcommand{\Signal}{X}
\newcommand{\Data}{Y}
\newcommand{\IMFs}{S}
\newcommand{\IMFt}{\widetilde S}
\newcommand{\IF}{F}
\newcommand{\IFP}{F_{\circ}}
\newcommand{\DFN}{D}
\newcommand{\fs}{f}
\newcommand{\fp}{f_{\circ}}
\newcommand{\xs}{x^{*}}
\newcommand{\xp}{x_{\circ}}
\newcommand{\xx}{x}

\newtheorem{proposition}{Proposition}[section]
\newtheorem{lemma}[proposition]{Lemma}
\newtheorem*{remark}{Remark}

\title{Proof of Lemma 4.2 and completion of Proposition 4.1}
\author{}
\date{}

\begin{document}
\maketitle

\setcounter{section}{3}
\section{Theoretical guarantees}

Let \(\mathcal T_n\subset[0,1]\) be the finite set of points at which the filter is
evaluated and put \(m=|\mathcal T_n|\).  For \(t\in\mathcal T_n\), write
\begin{equation}
  w_{u,t}\eqdef \Kern_h(t-u),\qquad
  W_t\eqdef\sum_u w_{u,t},\qquad
  Q_t^2\eqdef\sum_u w_{u,t}^2,\qquad
  k_t\eqdef\max_u w_{u,t}.
  \label{eq:weights}
\end{equation}
We assume \(w_{u,t}\geq0\), as for the Epanechnikov kernel used in
\texttt{IMF.tex}.  Suppose that \(\Data_u=\Signal_u+\varepsilon_u\), where the
\(\varepsilon_u\) are independent, and let
\begin{equation}
  \rhohum(z)\eqdef\E\rhohu(z+\varepsilon).
  \label{eq:phi}
\end{equation}
The argument below also allows independent, non-identically distributed errors if
\(\rhohum(\Signal_u-x)\) in each summand is replaced by
\(\E\rhohu(\Data_u-x)\).

Assume that \(\rhohu\in C^3(\R)\) is convex and, for finite constants
\(L_1,L_2,L_3\),
\begin{equation}
  |\rhohu'(z)|\leq L_1,\qquad
  0\leq\rhohu''(z)\leq L_2,\qquad
  |\rhohu'''(z)|\leq L_3,
  \qquad z\in\R.
  \label{eq:derivative-bounds}
\end{equation}
For the Gaussian-smoothed absolute-value contrast \(\rhohu_H\) in
\texttt{IMF.tex}, one may take
\begin{equation}
  L_1=1,\qquad
  L_2=\frac{2}{H\sqrt{2\pi}},\qquad
  L_3=\frac{2}{H^2\sqrt{2\pi}}.
  \label{eq:rhoH-constants}
\end{equation}

Define, with the notation of \texttt{IMF.tex},
\begin{align}
  \IMFt_t
  &\eqdef \arginf_{z\in\R}\sum_u\rhohu(\Data_u-z)w_{u,t},
  &
  \IMFs_t
  &\eqdef \arginf_{z\in\R}\sum_u\E\rhohu(\Data_u-z)w_{u,t},
  \label{eq:estimators}\\
  \IF_t
  &\eqdef\sum_u\rhohum''(\Signal_u-\IMFs_t)w_{u,t},
  &
  \DFN_t
  &\eqdef\IF_t^{1/2},
  \label{eq:curvature}\\
  \nabla_t
  &\eqdef\sum_u\left\{
       \rhohu'(\Data_u-\IMFs_t)
       -\E\rhohu'(\Data_u-\IMFs_t)
     \right\}w_{u,t}.
  \label{eq:nabla}
\end{align}
We require \(\IF_t>0\) for every \(t\in\mathcal T_n\).  This holds for
\(\rhohu_H\) whenever at least one \(w_{u,t}\) is positive.

For \(\xx\geq0\) and \(z_{\xx}\eqdef\xx+\log(4m)\), set
\begin{align}
  C_1
  &\eqdef
    \max_{t\in\mathcal T_n}
    \frac{L_1}{\DFN_t}
    \left\{Q_t\sqrt{2z_{\xx}}+\frac{2k_tz_{\xx}}{3}\right\},
  \label{eq:C1}\\
  C_2
  &\eqdef
    \max_{t\in\mathcal T_n}
    \frac{L_2}{\DFN_t}
    \left\{Q_t\sqrt{2z_{\xx}}+\frac{k_tz_{\xx}}{3}\right\},
  \label{eq:C2}\\
  C_3
  &\eqdef
    \max_{t\in\mathcal T_n}\frac{L_3W_t}{\DFN_t^2},
  &
  \DFN_{\min}
  &\eqdef\min_{t\in\mathcal T_n}\DFN_t,
  &
  \delta
  &\eqdef\frac{C_2}{\DFN_{\min}}.
  \label{eq:C3-delta}
\end{align}

The following statement supplies the missing right-hand side in Proposition~4.1.

\begin{proposition}\label{prop:main}
Suppose that the conditions above hold, that each empirical objective in
\eqref{eq:estimators} is strongly convex in the sense required by
Theorem~A.9 of \texttt{IRMF.tex}, and that
\begin{equation}
  \delta<1,
  \qquad
  \frac{C_1C_3}{\DFN_{\min}(1-\delta)^2}<\frac49.
  \label{eq:smallness}
\end{equation}
Then there is a random set \(\Omega(\xx)\) with
\(\Pp\{\Omega(\xx)\}\geq1-e^{-\xx}\) such that
\begin{equation}
  \sup_{t\in\mathcal T_n}
  \left|
    \DFN_t\bigl(
      \IMFt_t-\IMFs_t-\IF_t^{-1}\nabla_t
    \bigr)
  \right|
  \leq
  \frac{1}{\DFN_{\min}}
  \left\{
    \frac{3C_1^2C_3}{4(1-\delta)^3}
    +\frac{C_1C_2}{1-\delta}
  \right\}.
  \label{eq:proposition-bound}
\end{equation}
Consequently, the oracle-corrected estimator
\begin{equation}
  \IMFt_t^{\mathrm{corr}}
  \eqdef \IMFt_t-\IF_t^{-1}\nabla_t
  \label{eq:corrected-estimator}
\end{equation}
satisfies the same upper bound with
\(\IMFt_t-\IMFs_t-\IF_t^{-1}\nabla_t
=\IMFt_t^{\mathrm{corr}}-\IMFs_t\).
\end{proposition}

\begin{remark}[Sign of the correction]
With \(\nabla_t\) defined by \eqref{eq:nabla}, the minus sign in
\eqref{eq:proposition-bound} is forced by differentiation:
\(f_\circ'(\IMFs_t)=-\nabla_t\).  Thus
\(\IMFt_t-\IMFs_t\) has first-order term \(\IF_t^{-1}\nabla_t\), and one
subtracts that term to correct \(\IMFt_t\).  The plus sign in the unfinished
display in \texttt{IMF.tex} would be correct only if \(\nabla_t\) were defined
with the opposite sign.
\end{remark}

\subsection{Perturbation theory in action}

Fix \(t\in\mathcal T_n\).  All quantities in this subsection may depend on
\(t\).  Define
\begin{equation}
  \fs(z)\eqdef\sum_u\rhohum(\Signal_u-z)w_{u,t},
  \qquad
  \fp(z)\eqdef\sum_u\rhohu(\Data_u-z)w_{u,t},
  \label{eq:objectives}
\end{equation}
and
\begin{equation}
  \xs\eqdef\arginf_z\fs(z)=\IMFs_t,
  \qquad
  \xp\eqdef\arginf_z\fp(z)=\IMFt_t.
  \label{eq:minimizers}
\end{equation}
Further, put
\begin{equation}
  \IFP\eqdef\fp''(\xs),
  \qquad
  \IF\eqdef\fs''(\xs)=\DFN^2.
  \label{eq:F-Fcirc}
\end{equation}
For \(r>0\), let
\(U_t(r)\eqdef\{z\in\R:|\DFN(z-\xs)|\leq r\}\).

\begin{lemma}\label{lem:local}
On \(\Omega(\xx)\), simultaneously for every \(t\in\mathcal T_n\),
\begin{enumerate}
\item
\begin{align}
  \frac{1}{\DFN}
  |\fs'(\xs)-\fp'(\xs)|
  &=\frac{1}{\DFN}
    \left|\sum_u
      \left\{\rhohu'(\Data_u-\xs)
      -\E\rhohu'(\Data_u-\xs)\right\}w_{u,t}
    \right|
  \leq C_1,
  \label{eq:lemma-score}\\
  \frac{1}{\DFN}
  |\fs''(\xs)-\fp''(\xs)|
  &=\frac{1}{\DFN}
    \left|\sum_u
      \left\{\rhohu''(\Data_u-\xs)
      -\E\rhohu''(\Data_u-\xs)\right\}w_{u,t}
    \right|
  \leq C_2.
  \label{eq:lemma-hessian}
\end{align}

\item
\begin{equation}
  \sup_{z\in\R}
  \frac{1}{\DFN^2}
  \left|\sum_u\rhohu'''(\Data_u-z)w_{u,t}\right|
  \leq C_3.
  \label{eq:lemma-third}
\end{equation}

\item For every \(r>0\),
\begin{equation}
  \sup_{z\in U_t(r)}
  \frac{1}{\DFN^2}
  \left|\sum_u
    \left\{
      \rhohum''(\Signal_u-z)
      -\rhohum''(\Signal_u-\xs)
    \right\}w_{u,t}
  \right|
  \leq \frac{C_3r}{\DFN}.
  \label{eq:lemma-population-hessian}
\end{equation}
\end{enumerate}
\end{lemma}

\begin{proof}
Differentiating \eqref{eq:objectives} gives
\begin{align*}
  \fs'(z)&=-\sum_u\rhohum'(\Signal_u-z)w_{u,t},
  &
  \fp'(z)&=-\sum_u\rhohu'(\Data_u-z)w_{u,t},\\
  \fs''(z)&=\sum_u\rhohum''(\Signal_u-z)w_{u,t},
  &
  \fp''(z)&=\sum_u\rhohu''(\Data_u-z)w_{u,t}.
\end{align*}
Since differentiation under the expectation is justified by
\eqref{eq:derivative-bounds},
\(\rhohum^{(j)}(a)=\E\rhohu^{(j)}(a+\varepsilon)\) for
\(j=1,2,3\).  The identities in part~(i) follow.

For the first identity, set
\[
  Z_{u,t}\eqdef w_{u,t}
  \left\{\rhohu'(\Data_u-\xs)
  -\E\rhohu'(\Data_u-\xs)\right\}.
\]
These variables are independent and centered.  Moreover,
\[
  |Z_{u,t}|\leq2L_1k_t,
  \qquad
  \sum_u\E Z_{u,t}^2\leq L_1^2Q_t^2.
\]
Bernstein's inequality therefore yields, for every \(z>0\),
\begin{equation}
  \Pp\left(
    \left|\sum_uZ_{u,t}\right|>
    L_1\left\{Q_t\sqrt{2z}+\frac{2k_tz}{3}\right\}
  \right)
  \leq2e^{-z}.
  \label{eq:bernstein-score}
\end{equation}
For the second identity, use
\[
  \widetilde Z_{u,t}\eqdef w_{u,t}
  \left\{\rhohu''(\Data_u-\xs)
  -\E\rhohu''(\Data_u-\xs)\right\}.
\]
Convexity and \eqref{eq:derivative-bounds} give
\[
  |\widetilde Z_{u,t}|\leq L_2k_t,
  \qquad
  \sum_u\E\widetilde Z_{u,t}^2\leq L_2^2Q_t^2.
\]
Hence
\begin{equation}
  \Pp\left(
    \left|\sum_u\widetilde Z_{u,t}\right|>
    L_2\left\{Q_t\sqrt{2z}+\frac{k_tz}{3}\right\}
  \right)
  \leq2e^{-z}.
  \label{eq:bernstein-hessian}
\end{equation}
Take \(z=z_{\xx}=\xx+\log(4m)\) in
\eqref{eq:bernstein-score} and \eqref{eq:bernstein-hessian}, and take the
union over the two inequalities and over \(t\in\mathcal T_n\).  The exceptional
probability is at most
\(4m e^{-z_{\xx}}=e^{-\xx}\).  Equations \eqref{eq:C1} and
\eqref{eq:C2} now prove part~(i).

Part~(ii) is deterministic.  By \eqref{eq:derivative-bounds} and
\eqref{eq:C3-delta},
\[
  \sup_{z\in\R}\frac{1}{\DFN^2}
  \left|\sum_u\rhohu'''(\Data_u-z)w_{u,t}\right|
  \leq\frac{L_3W_t}{\DFN^2}\leq C_3.
\]
Finally, \(|\rhohum'''|\leq L_3\).  The mean value theorem and
\(|z-\xs|\leq r/\DFN\) for \(z\in U_t(r)\) give
\begin{align*}
  &\frac{1}{\DFN^2}
  \left|\sum_u
    \left\{\rhohum''(\Signal_u-z)
    -\rhohum''(\Signal_u-\xs)\right\}w_{u,t}\right|\\
  &\hspace{35mm}\leq
    \frac{L_3W_t}{\DFN^2}|z-\xs|
    \leq\frac{C_3r}{\DFN},
\end{align*}
which proves part~(iii).
\end{proof}

\begin{proof}[Proof of Proposition~\ref{prop:main}]
Fix \(t\in\mathcal T_n\).  Since \(\fs'(\xs)=0\),
\begin{equation}
  \fp'(\xs)
  =-\sum_u\left\{
      \rhohu'(\Data_u-\xs)-\E\rhohu'(\Data_u-\xs)
    \right\}w_{u,t}
  =-\nabla_t.
  \label{eq:score-sign}
\end{equation}
Lemma~\ref{lem:local}(i) also gives
\begin{equation}
  \frac{|\IFP-\IF|}{\IF}
  \leq\frac{C_2}{\DFN}
  \leq\delta.
  \label{eq:F-comparison}
\end{equation}
In particular, \(\IFP\geq(1-\delta)\IF>0\).  With
\begin{equation}
  b_{\circ}\eqdef
  \frac{\DFN}{\IFP}|\fp'(\xs)|,
  \label{eq:bcirc}
\end{equation}
the score bound and \eqref{eq:F-comparison} imply
\begin{equation}
  b_{\circ}\leq\frac{C_1}{1-\delta}.
  \label{eq:bcirc-bound}
\end{equation}

Lemma~\ref{lem:local}(ii) gives the third-derivative parameter
\(\tau_3=C_3/\DFN\) in the local metric \(\DFN\).  Also,
\eqref{eq:F-comparison} permits
\(\kappa^2=(1-\delta)^{-1}\) in the comparison
\(\DFN^2\leq\kappa^2\IFP\).  Thus the second condition in
\eqref{eq:smallness} implies
\begin{equation}
  \kappa^2\tau_3 b_{\circ}
  \leq\frac{C_1C_3}{\DFN(1-\delta)^2}<\frac49.
\end{equation}
The generic perturbation bound, Theorem~A.9 in \texttt{IRMF.tex}, now gives
\begin{equation}
  \left|\DFN^{-1}\IFP
    \left(\xp-\xs+\IFP^{-1}\fp'(\xs)\right)\right|
  \leq\frac{3C_3}{4\DFN}b_{\circ}^2.
  \label{eq:generic-bound}
\end{equation}
The third-derivative bound is global, so its neighborhood condition may be
satisfied with any radius at least \(3b_{\circ}/2\).  Combining
\eqref{eq:F-comparison}, \eqref{eq:bcirc-bound}, and
\eqref{eq:generic-bound} yields
\begin{equation}
  \left|\xp-\xs+\IFP^{-1}\fp'(\xs)\right|
  \leq
  \frac{3C_1^2C_3}{4\DFN^2(1-\delta)^3}.
  \label{eq:Fcirc-remainder}
\end{equation}

It remains to replace \(\IFP\) by the population curvature \(\IF\).  By
\eqref{eq:F-comparison}, \eqref{eq:lemma-score}, and
\(\fs'(\xs)=0\),
\begin{align}
  \left|\IFP^{-1}\fp'(\xs)-\IF^{-1}\fp'(\xs)\right|
  &=\frac{|\IFP-\IF|}{\IFP\IF}|\fp'(\xs)|
  \notag\\
  &\leq\frac{C_1C_2}{\DFN^2(1-\delta)}.
  \label{eq:F-replacement}
\end{align}
The triangle inequality applied to \eqref{eq:Fcirc-remainder} and
\eqref{eq:F-replacement}, followed by \eqref{eq:score-sign}, proves
\begin{equation}
  \DFN\left|
    \xp-\xs-\IF^{-1}\nabla_t
  \right|
  \leq
  \frac{1}{\DFN}
  \left\{
    \frac{3C_1^2C_3}{4(1-\delta)^3}
    +\frac{C_1C_2}{1-\delta}
  \right\}.
  \label{eq:pointwise-final}
\end{equation}
Taking the supremum over \(t\in\mathcal T_n\) and using
\(\DFN_t\geq\DFN_{\min}\) gives \eqref{eq:proposition-bound}.
\end{proof}

\begin{remark}[What the source files do not establish]
The union bound above proves the result on the finite IMF evaluation grid.  A
literal supremum over every \(t\in[0,1]\) needs a covering-number argument and
regularity bounds in \(t\) for \(\IMFs_t\), \(\IF_t\), and the kernel-weighted
process.  Those assumptions are not present in \texttt{IMF.tex} or
\texttt{IRMF.tex}.  They also do not state the uniform positive-curvature and
smallness assumptions in \eqref{eq:smallness}, which Theorem~A.9 needs.  Finally,
the source files define \(N_t=N_t(\Signal_t)\), whereas
\(\IF_t=N_t(\IMFs_t)\).  Hence the sentence ``\(\IF=N_t\)'' is valid only if
\(N_t\) is redefined at \(\IMFs_t\), or if \(\Signal_t=\IMFs_t\).
\end{remark}

\end{document}
